\chapter{Coding Sheet}\label{coding_sheet}
This appendix contains the coding sheet used for the deductive content analysis of the interviews. The coding sheet is based on the taxonomy of misconceptions presented in \autoref{taxonomy}, and contains definitions and examples for each category in the taxonomy. The coding sheet was used by the coders to identify and classify misconceptions in the interview transcripts.

\section*{Dimension 1: Foundational Conceptual Accuracy}

\subsection*{Anthropomorphism}

\textit{Attributing human mental states, intentions or understanding to a LLM.}

This is a quite broad category, but subjects should be placed in this category if they in some way expressed belief in attributes not fitting of a statistical model and such attributes are aligned with human mental states.

\subsection*{Deterministic Automation}

\textit{Confusing LLMs with deterministic automation systems.}

This could be deterministic chatbot systems, but also more traditional automation systems such as a calculator, or algorithms in general. Subjects who in some way expressed belief that LLMs operate based on deterministic rules, or that they are deterministic systems should be categorized into this category.

\subsection*{Accuracy Overestimation}

\textit{Overestimating the reliability and correctness of LLM outputs.}

Subjects who in some way expressed belief that LLMs are more reliable, correct or trustworthy than they actually are should be categorized into this category.

\subsection*{Autonomy Overestimation}

\textit{Belief that the LLM operate independently or possess agency.}

Subjects who in some way expressed belief that LLMs operate independently, or possess agency should be categorized into this category. This could be belief in some sort of action from the model other than generating output, like ``searching the internet''.

\section*{Dimension 2: Mechanistic Understanding}

\subsection*{Real-time Data Belief}

\textit{Belief that the LLM retrieves current information.}

Subjects who in some way expressed belief that LLMs retrieve current information should be categorized into this category. This retrieval can be from the internet, from some database, or from the prompt itself.

\subsection*{Hallucination Blindness}

\textit{Failure to recognize that LLMs can produce false or fabricated information.}

This is not just about recognizing that LLMs can produce false information, but also recognizing that hallucinations is part of the core idea of how LLMs work, and that it is not just a bug or an error that can be fixed.

\subsection*{Context Window Neglect}

\textit{Failure to recognize limits of LLMs' context window.}

If subjects shows no indication of being aware of the limits of LLMs' context window should be categorized into this category. This could be expressed as belief that the model can keep track of an entire conversation, or that it can remember information from previous conversations.

\subsection*{Semantic Understanding}

\textit{Lack of awareness that the LLM processes language as tokens rather than semantic units.}

Subjects who in some way expressed belief that LLMs understand language on a semantic level, or that they read the text as given to them.

\subsection*{Source Attribution Error}

\textit{Belief that LLM output originate from specific identifiable sources.}

Subjects who in some way expressed belief that LLM output originate from specific identifiable sources should be categorized into this category.

\subsection*{Conditional Treatment}

\textit{Belief that LLMs process prompts differently based on the content or context.}

Subjects who in some way expressed belief that LLMs will process one prompt differently than another should be categorized into this category.

\section*{Dimension 3: Ethical and Academic Integrity}

\subsection*{Plagiarism Blindness}

\textit{Failure to recognize ethical and legal issues related to the traceability of sources behind LLM-generated content.}

This is strongly correlated to the source attribution error, but is more focused on the ethical and legal issues related to the issue. A LLM will mimmic citation patterns in the training data, and failure to recognize the ethical and legal issues related to the traceability of sources behind LLM-generated content should be categorized into this category.

\subsection*{Copyright Blindness}

\textit{Failure to recognize that LLM-generated content may still be subject to copyright and usage restrictions, despite being machine-generated.}

This category is focused on the subjects ability to recognize the questions surrounding ownership of LLM-generated content. That is wether it can be claimed by the prompt creator, the LLM creator, or the original creator of the content in the training data. Failure to recognize this issue should be categorized into this category.

\subsection*{Conditional Acceptance}

\textit{Ethical acceptance of LLM use only under certain conditions or domains.}

Subjects who in some way expressed that ``use of LLMs is fine if used for \ldots'' or ``LLM use is not fine in my own field, but in \ldots'' should be categorized into this category.

\subsection*{Environmental Neglect}

\textit{Failure to recognize environmental impact of LLMs.}

Subjects who in some way expressed belief that LLMs have negligible environmental impact, or showed lack of awareness of energy consumption associated with LLM training and operation should be categorized into this category.

\section*{Dimension 4: Socio-Cognitive Impact}

\subsection*{Replacement Fear}

\textit{Belief that LLMs will replace human roles or professions, expertise, or professions.}

Subjects who in some way expressed belief that LLMs will replace human roles or professions should be categorized into this category. Subjects should not be categorized into this category if they only expressed belief that LLMs will replace some tasks within a profession.

\subsection*{Cognitive Passivity}

\textit{Tendency to rely on LLMs rather than engaging in independent reasoning or critical thinking.}

Subjects who, when they described how they use LLMs, showed a tendency to rely on LLMs for tasks that would normally require independent reasoning or critical thinking, should be categorized into this category.

\subsection*{Dependency Risk}

\textit{Failure to recognize that extensive can lead to long term reliance on LLMs, skill degradation or loss of autonomy.}

Subjects who, in describing their use of LLMs, showed no recognition of the risk of becoming reliant on LLMs or described tasks that should not require LLMs but still used them for such tasks, should be categorized into this category.

\subsection*{Credibility Bias}

\textit{Treating LLM outputs as inherently truthful or authoritative without critical evaluation.}

Subjects could think of LLM output as reliable just because of the fact that is machine generated. If this is shown in the subjects description of their use of LLMs, they should be categorized into this category.

